% =====================================================================
\section{Словарь терминов}
% =====================================================================
{\small
\begin{xltabular}{\textwidth}{@{}>{\bfseries}p{40mm} X r@{}}
\toprule
\normalfont Термин & Что означает & Урок \\
\midrule
\endfirsthead
\toprule
\normalfont Термин & Что означает & Урок \\
\midrule
\endhead
1-Wire & шина с одним проводом данных и подтяжкой; на ней может работать несколько датчиков с уникальными адресами (например, DS18B20) & 15 \\
BLE & Bluetooth Low Energy: экономичный радиоканал для связи с телефоном рядом с устройством & 12 \\
Backtrace & цепочка адресов функций, которые выполнялись в момент сбоя & 14 \\
Brownout & аварийный сброс из-за просадки напряжения питания & 14 \\
CS (Chip Select) & линия SPI, которой ведущий выбирает одно ведомое устройство & 10 \\
Deep-sleep & глубокий сон: работает только RTC-домен, пробуждение --- это перезапуск программы & 9 \\
Duty cycle & коэффициент заполнения ШИМ: доля времени, когда сигнал в HIGH & 4 \\
DVP & параллельный интерфейс камеры: 8 линий данных и сигналы PCLK, VSYNC, HREF & 20 \\
ESP-NOW & протокол Espressif для коротких сообщений между платами без роутера & 12 \\
FreeRTOS & операционная система реального времени, на которой работает Arduino-ESP32 & 7 \\
Git & система хранения истории изменений проекта в виде снимков (коммитов) & 28 \\
GPIO & вывод общего назначения: цифровой вход или выход микроконтроллера & 0, 1 \\
H-мост & схема из четырёх ключей, позволяющая менять полярность на моторе, то есть направление вращения & 16 \\
HID & класс USB-устройств ввода: клавиатура, мышь, мультимедийные клавиши & 19 \\
HSV & цветовая модель: оттенок, насыщенность, яркость; удобна для радуги & 17 \\
I\textsuperscript{2}C & двухпроводная шина с адресами устройств: SDA (данные) и SCL (такт) & 6 \\
IMU & инерциальный датчик: акселерометр и гироскоп в одном корпусе (например, MPU6050) & 26 \\
ISR & обработчик прерывания: короткая функция, которую процессор вызывает по событию & 3 \\
JPEG & формат сжатого изображения; камера OV2640 формирует его сама & 20 \\
JSON & текстовый формат данных: пары «имя: значение» в фигурных скобках & 18 \\
LEDC & периферия ESP32-S3 для формирования ШИМ-сигналов & 4 \\
Li-ion & литий-ионный аккумулятор: \SIrange{3.0}{4.2}{V}, требует платы защиты и специальной зарядки & 27 \\
LittleFS & файловая система во Flash: файлы и папки без SD-карты & 18 \\
MJPEG & видеопоток из последовательности JPEG-кадров & 20 \\
MOSFET & полевой транзистор, управляемый напряжением на затворе; ключ для мощной нагрузки & 16 \\
MQTT & протокол обмена сообщениями через брокер по темам (topics) & 12 \\
NTP & протокол синхронизации времени с серверами в сети & 18 \\
NVS & хранилище «ключ--значение» во Flash для настроек & 11 \\
OTA & обновление прошивки по сети без USB-кабеля & 13 \\
PSRAM & внешняя оперативная память модуля для больших буферов & 11 \\
RMT & периферия ESP32-S3 для точных последовательностей импульсов (WS2812, ИК-пульты) & 17 \\
RTC-память & небольшая память, которая сохраняется в deep-sleep & 9 \\
SCCB & двухпроводная шина настройки камеры, совместимая с I\textsuperscript{2}C & 20 \\
SPI & быстрая синхронная шина: SCK, MOSI, MISO и CS & 10 \\
Strapping-выводы & выводы, уровень которых читается при сбросе и задаёт режим загрузки & 3 \\
TLS & протокол шифрования и проверки подлинности сервера (основа HTTPS) & 13 \\
UART & асинхронный последовательный порт: TX, RX и общая земля & 2 \\
Watchdog & сторожевой таймер: перезагружает чип, если программа зависла & 14 \\
WS2812 & адресный RGB-светодиод с микросхемой внутри; управляется одним проводом & 17 \\
Адресный светодиод & светодиод с микросхемой, цвет которого задаётся данными по одному проводу & 17 \\
Акселерометр & датчик ускорения; по направлению силы тяжести показывает наклон & 26 \\
АЦП & аналого-цифровой преобразователь: превращает напряжение в число & 5 \\
Аттенюатор & делитель на входе АЦП, расширяющий диапазон измеряемых напряжений & 5 \\
Бит, байт & бит --- один двоичный разряд (0 или 1); байт --- восемь битов & 22 \\
Гироскоп & датчик скорости поворота вокруг осей & 26 \\
Делитель напряжения & два резистора, на среднем выводе которых получается часть входного напряжения & 5, 15 \\
Драйвер мотора & микросхема с H-мостом, которая управляет мотором по сигналам платы & 16 \\
Дребезг & многократное замыкание-размыкание механического контакта при нажатии & 3 \\
Задача & независимая функция с бесконечным циклом под управлением FreeRTOS & 7 \\
Класс & тип данных, объединяющий переменные и функции для работы с ними & 28 \\
Комплементарный фильтр & способ смешать показания гироскопа и акселерометра, убирая шум и дрейф & 26 \\
Конечный автомат & описание поведения устройства через состояния, события и переходы & 21 \\
Маска & число с единицами в нужных битах; используется в битовых операциях & 22 \\
Мьютекс & «ключ», который защищает общие данные от одновременного доступа задач & 7 \\
Оптопара & прибор, передающий сигнал светом и электрически разделяющий цепи & 16 \\
Очередь & потокобезопасный буфер для передачи данных между задачами & 7 \\
Подтяжка & резистор, задающий уровень на входе, пока к нему ничего не подключено & 3 \\
Прерывание & сигнал, по которому процессор откладывает основную программу и вызывает ISR & 3 \\
Регистр & ячейка памяти, биты которой управляют аппаратурой чипа или датчика & 22, 26 \\
Реле & электромагнитный выключатель, механически замыкающий контакты & 16 \\
Сдвиговый регистр & микросхема, принимающая биты по одному проводу и выдающая их на несколько выходов (74HC595) & 22 \\
Тайм-аут & предел времени ожидания, после которого программа перестаёт ждать & 15, 18 \\
Таймер аппаратный & счётчик внутри чипа, вызывающий прерывание через точные интервалы & 21 \\
Шаговый двигатель & двигатель, поворачивающийся одинаковыми шагами по командам программы & 25 \\
ШИМ & широтно-импульсная модуляция: управление средним уровнем долей времени включения & 4 \\
Энкодер & ручка без упоров, сообщающая направление и число щелчков поворота & 23 \\
I²S & цифровой поток аудио с линиями BCLK, WS и данных & 29 \\
RMT & аппаратная передача и приём последовательностей уровней заданной длительности & 30 \\
PCNT & аппаратный счётчик фронтов & 31 \\
MCPWM & периферия управления связанными PWM-выходами & 31 \\
TWAI & контроллер Classical CAN; нужен внешний трансивер & 32 \\
DMA & обмен блоками между памятью и периферией без обслуживания каждого байта CPU & 29, 33 \\
PHY & физический приёмопередатчик интерфейса & 34 \\
JTAG & интерфейс отладки процессора & 35 \\
ULP & маломощный сопроцессор с отдельной программой & 36 \\
Инференс & применение обученной модели к новому входу & 37 \\
Provisioning & первичная выдача устройству настроек и учётных данных & 38 \\
WebSocket & постоянный двусторонний канал между браузером и сервером & 39 \\
Callback & функция, которую библиотека вызывает при событии & 39 \\
ESPUI & конструктор веб-панели из C++ с готовыми элементами & 39 \\
Captive portal & страница входа или настройки, на которую направляют подключённый браузер & 40 \\
WiFiManager & библиотека выдачи устройству настроек Wi-Fi через портал & 40 \\
\bottomrule
\end{xltabular}
}
